\begin{textsample}{POS Dim 3 – persona\_grok – Score 17.00 – t572\_grok.txt}  \label{ex:f3_pos_038}
In our modern world, a culture of extravagant \textbf{consumption} is fueling an alarming surge in electronic \textbf{waste}, with devastating consequences for the planet.
Take devices like iPhones, for instance, where planned obsolescence is built right in. \textbf{Companies} such as Apple create \textbf{products} that are notoriously hard to repair and actively lobby against right-to-repair laws, ensuring that gadgets become obsolete far too quickly.
This \textbf{cycle} begins with destructive mining for essential \textbf{minerals}, which ravages ecosystems and communities.
Devices have shockingly short lifespans, contributing to a global e-waste \textbf{problem} that now exceeds 65 million \textbf{tonnes} annually.
That 's a mountain of discarded tech poisoning landfills and oceans.
Our latest Greenpeace \textbf{report} puts the spotlight on corporate responsibility by grading major firms on their repairability efforts.
We 're praising \textbf{companies} like HP, Dell, and Fairphone for leading the way in making devices easier to fix and extend.
On the other hand, we 're calling out Samsung and others for falling short and perpetuating the throwaway \textbf{model}.
To turn this around, we \textbf{need} a fundamental \textbf{shift} to circular \textbf{production} methods that prioritize \textbf{reuse} and \textbf{recycling}.
Strong policies promoting \textbf{product} durability are essential, alongside a move toward modest \textbf{consumption}.
These changes can help curb ecological overshoot and bolster the \textbf{growth} of renewable \textbf{energy}, paving the way for a more sustainable future.

% matched lemmas: company, consumption, cycle, energy, growth, mineral, model, need, problem, product, production, recycling, report, reuse, shift, tonne, waste
\end{textsample}
